\documentclass[11pt,a4paper]{article}
\usepackage[UTF8]{ctex}
\usepackage{geometry}
\geometry{left=2.5cm,right=2.5cm,top=2.5cm,bottom=2.5cm}
\usepackage{amsmath,amssymb,amsthm}
\usepackage{hyperref}
\usepackage{booktabs}
\usepackage{enumitem}
\usepackage{xltabular}
\usepackage{array}
\hypersetup{colorlinks=true,linkcolor=blue,urlcolor=blue}

\newtheorem{definition}{定义}[section]
\newtheorem{theorem}{定理}[section]
\newtheorem{proposition}{命题}[section]
\newtheorem{remark}{注记}[section]

\title{安全运营最优调度理论的学术定位与对话\\
——对抗性博弈与成本敏感决策的对比分析}
\author{李龙胜}
\date{\today}

\begin{document}

\maketitle

\begin{center}
\textit{
当前学术界在告警优先级排序领域已积累大量研究成果，\\
但所有研究都在追求“更好”——更好的分类器、更优的调度算法。\\
我们追问一个被忽视的问题：多好才是够好？\\
本文将贪心择优+参数扰动与对抗强化学习进行系统对比，\\
将够用就行决策与成本敏感决策进行深度对接，\\
确立理论在学术版图中的独特定位。\\
}
\end{center}

\vspace{5mm}

\begin{abstract}
安全运营最优调度理论的核心主张是：安全运营中的所有决策，
都是在有限资源下做最优分配，而最优的判据是边际收益等于边际成本。
本文将该理论与学术界两个最相关的研究方向进行系统对话。
在对抗性博弈方向上，与AAAI 2023的对抗强化学习方法进行对比，
论证贪心择优+参数扰动在计算效率、理论简洁性和可解释性上的优势。
在成本敏感决策方向上，与2026年Trust Alignment方法进行对接，
展示理论如何为成本敏感决策提供从公设到落地工具的完整实例化。
在此基础上，明确理论在学术版图中的统一最优性判据定位，
为后续的学术发表和理论推广提供文献基础。
本文基于厂商调查报告和学术文献综述两份前置文档的发现，
进一步深化了理论与工程实践和学术前沿的对话。
\end{abstract}

\tableofcontents

\section{引言：理论与学术前沿的对话需求}

安全运营最优调度理论已经建立了完整的理论框架——
从公设推导贪心择优定理，从台账量化经济损失，
从柯克霍夫原则保护核心参数，从五表四步模型实现工程落地。
然而，理论不能孤立存在。
它需要与学术界已有的研究成果进行对话，
以确立自身的学术定位，识别真正的差异化贡献。

在已完成的两份前置文档——
《安全厂商告警降噪技术方案调查报告》和
《告警优先级排序与降噪学术研究个人梳理》——
的基础上，本文选取学术界与理论最直接相关的两个研究方向，
进行系统性的对比分析。

\section{对抗性博弈：贪心择优 vs 对抗强化学习}

\subsection{AAAI 2023对抗RL方法的核心设定}

AAAI 2023发表的《在移动的干草堆中找针：用对抗强化学习进行告警优先级排序》
是学术界与理论最直接相关的论文。
其核心设定与理论的柯克霍夫原则在逻辑上完全同构：
攻击者完全知道防御方的检测系统状态和告警优先级策略，
会动态选择最优攻击方式。

该研究提出用对抗强化学习计算鲁棒的告警优先级策略。
具体方法结合了双预言框架与神经强化学习——
防御方和攻击方的预言者分别用强化学习计算
对对方混合策略的最优响应。
在每一轮迭代中，
攻击方的RL agent学习如何绕过当前的防御策略，
防御方的RL agent学习如何应对攻击方的最新策略。
通过多轮迭代，系统收敛到攻防双方的一个均衡策略对。

\subsection{对抗RL方法的局限性}

然而，对抗RL方法面临几个根本性的局限。

\textbf{计算复杂度极高}：
每一轮迭代需要训练两个RL模型，
且双预言的收敛需要足够的迭代次数。
在真实SOC环境中，告警类型数量巨大、
攻防策略空间高维，对抗RL的部署成本可能难以承受。

\textbf{依赖训练数据}：
RL模型的训练需要大量攻防交互数据，
而真实攻防数据稀缺且高度敏感，
大部分研究只能在仿真环境中验证。

\textbf{缺乏对最优策略的解析刻画}：
对抗RL的输出是一个数值策略，
无法回答《\textit{为什么}这个策略是最优的》、
《\textit{在什么条件下}这个策略成立》等问题。
它提供的是数值近似，而非解析证明。

\subsection{贪心择优+参数扰动的替代方案}

理论提供了一种计算上更高效、理论上更简洁的替代方案。

\textbf{贪心择优定理保证最优性}：
在进展单调性和成本凸性的条件下，
每步贪心选择等价于全局最优。
这一结论不是通过数值逼近得到的，
而是从公设出发严格推导出来的。

\textbf{参数扰动对抗探测}：
通过周期性对台账参数施加微小扰动，
使攻击方在上个周期积累的探测知识强制过期。
参数扰动的有效性有信息论下界的严格证明——
攻击方要精确估计防御方参数，
需要数万次探测，远超单周期探测预算。

\textbf{计算效率的显著优势}：
贪心择优不需要训练任何模型，只需比较成本大小。
参数扰动不需要对抗训练，只需对台账施加随机噪声。
整个方案可以在现有SOC硬件上部署运行。

\subsection{理论对话总结}

对抗RL和贪心择优+参数扰动，
代表了两条不同的技术路线。
对抗RL通过数值逼近寻找最优策略，
贪心择优通过解析证明推导最优策略。
前者依赖训练数据和计算资源，
后者只依赖成本函数的凸性和进展量的单调性。
在工程上，贪心择优更轻量、更可解释、更易部署。
在理论上，贪心择优更优雅——
它不依赖任何外部假设，直接从生成论公设中推导出来。

\section{成本敏感决策：够用就行 vs Trust Alignment}

\subsection{2026年Trust Alignment方法的核心贡献}

2026年发表的《决策感知的信任信号对齐 for SOC告警分诊》
明确提出了成本敏感的概念：
漏报比误报昂贵得多，
需要用校准后的置信度和成本敏感决策阈值
来构建决策支持层。

该研究在经济学视角上与理论最接近。
它已经意识到告警优先级排序不是简单的分类问题，
而是成本-收益权衡问题。
它提出了《决策感知》的框架，
让AI系统的输出不仅考虑准确率，还考虑不同类型错误的成本差异。

\subsection{Trust Alignment方法未回答的问题}

然而，Trust Alignment方法只提出了《应该成本敏感》的方向，
没有给出完整的理论框架。
它未能回答以下几个核心问题。

\textbf{成本从哪来}：
不同类型的错误成本如何量化？
如何从企业业务数据中映射出具体的损失值？

\textbf{如何在不同告警类型之间分配分析资源}：
在总分析能力有限的前提下，
哪些告警应该优先分析，哪些可以忽略？
分配的数学判据是什么？

\textbf{如何处理攻击方的自适应行为}：
攻击方会主动适应防御策略的成本设置，
如何防止攻击方通过探测来操纵成本参数？

\textbf{如何将安全决策从《分类》升级为《资源分配》}：
告警优先级排序不只是《找得更准》，
而是在有限资源下做最优分配。
两者的数学结构有何不同？

\subsection{够用就行理论的完整回答}

理论恰好完整回答了这些问题。

\textbf{成本从台账中来}：
A值是漏报损失（从业务资产价值中映射），
\(\kappa\) 是分析成本（从工时和时薪中计算），
C是总分析能力预算（从团队规模和有效工时中估计）。

\textbf{量化通过标定和校准机制}：
从拍脑袋开始，通过反馈校准（漏报翻倍、长期误报衰减）
和全量回溯（季度重新评估所有告警类型），
逐步逼近真实的A值分布。

\textbf{最优分配通过贪心择优}：
按边际收益降序排列告警类别，
从低价值告警开始削减分析率，
直到总工时消耗降至预算以内。
最优分析率 \(r^*\) 的判据是
\(\alpha A > \kappa\) 时分析，否则忽略。

\textbf{对抗性通过柯克霍夫原则}：
参数保密——台账参数加密存储，厂商和攻击方均不可访问；
周期性扰动——每次演练后微调参数，
使攻击方探测知识强制过期；
伪造响应——以一定概率伪造分析/忽略结果，
让攻击方面临欠定推断困境。

\textbf{资源分配通过四步调度模型}：
来源过滤修正A值、
最优分析率计算、
设备分流确定分析路径、
配额动态调整与反馈校准。

\subsection{理论对话总结}

Trust Alignment方法识别了成本敏感决策的重要性，
够用就行理论为成本敏感决策提供了
从公设到落地工具的完整实例化。
前者指出了方向，后者完成了全链条推导。
前者提出了《应该考虑成本》的倡议，
后者给出了《如何量化成本、如何分配资源、
如何对抗操纵》的可操作方案。

\section{理论的学术定位}

基于以上两个方向的系统对比，
安全运营最优调度理论在学术版图中可以确立以下定位。

\textbf{在对抗性博弈方向上}：
提供了一种计算上更高效、理论上更简洁的替代方案。
贪心择优定理从公设层面严格证明了策略的最优性，
参数周期性扰动有信息论下界的严格支撑。
这两者共同构成了对抗性安全运营的统一理论框架。

\textbf{在成本敏感决策方向上}：
提供了从公设到落地的完整理论体系。
Trust Alignment等方法指出了成本敏感的方向，
够用就行理论完成了这一方向的完整实例化——
从成本量化到资源分配到对抗防御到工程落地。

\textbf{在告警优先级排序领域整体}：
理论最核心的差异化优势是提供了一个统一的最优性判据——
边际收益等于边际成本。
当前所有学术研究都在追求《更好》——
更好的分类器、更优的调度算法、更高的降噪率。
但没有人追问《多好才是够好》。
理论填补的正是这个空白。

\section{诚实标注}

\begin{center}
\begin{xltabular}{\textwidth}{c|X|c}
\toprule
\textbf{命题} & \textbf{状态} & \textbf{层级} \\
\midrule
对抗性博弈对比分析 &
与AAAI 2023论文的问题设定同构，
贪心择优在计算效率和理论简洁性上的优势已论证 &
II（系统性对比已完成） \\
成本敏感决策对比分析 &
与Trust Alignment论文的经济学视角对接，
够用就行理论的完整度优势已论证 &
II（系统性对比已完成） \\
统一最优性判据定位 &
《边际收益=边际成本》作为领域统一判据的论证已展开，
当前所有研究均未提出此判据 &
II（理论独创性已确立） \\
厂商降噪率竞赛现象 &
厂商调查报告已实证验证，
降噪率竞赛确实存在且缺乏最优性标准 &
II（实证支持已到位） \\
\bottomrule
\end{xltabular}
\end{center}

\section{结语}

安全运营最优调度理论不再是孤立的理论推演。
它在对抗性博弈方向上找到了与学术前沿的精确对话点，
在成本敏感决策方向上找到了经济学视角的学术同盟。
理论最核心的差异化优势——
为整个领域提供统一的最优性判据——
在此得到了学术文献和工程实践的双重支撑。

接下来，是向学术界和对理论真正感兴趣的企业展示，
让他们看到我们在思路上独有的东西。
理论的纸面工作已如此坚实。

\vspace{5mm}
\noindent\textbf{文档信息}：
\begin{itemize}
    \item 作者：李龙胜
    \item 版本：第一版（学术定位与对话归档）
    \item 许可：CC BY 4.0
    \item 前置依赖：四卷体系、安全厂商调查报告、
          告警优先级排序与降噪学术研究个人梳理、攻防博弈论统一框架
\end{itemize}

\end{document}